\section{Background: MCP \& A2A}

Modern agents no longer interact with a small, fixed set of tools, rather they must also communicate with peer agents to complete complex tasks successfully. The resulting multi-agent systems enable the decomposition of complex problems among specialized agents, each contributing distinct capabilities and domain knowledge within a collaborative network~\cite{krishnan2025advancingmas}. Two protocols currently dominate this landscape: MCP for tool invocation and A2A for agent-to-agent communication.

\subsection{Model Context Protocol}
MCP is an open protocol that standardizes the interface through which large language model applications discover and invoke external capabilities~\cite{mcp2024introduction}. It adopts a client-server architecture in which an MCP host connects through lightweight clients to one or more MCP servers, each exposing a set of \emph{tools}, \emph{resources}, and \emph{prompts}. A tool definition comprises a unique name, a natural-language description intended to inform the model's selection process, and a JSON schema specifying the expected input parameters. At runtime, the full set of available tool definitions is typically injected into the LLM's context window, providing the information necessary for the model to determine which tool to invoke and how to construct a valid call. MCP thereby decouples capability access from backend-specific implementation details, yet delegates both discovery and execution reasoning to the consuming model within a single context.

\subsection{Agent-2-Agent Protocol}
A2A, on the other hand, addresses the discovery of and collaboration between autonomous agents~\cite{google2024a2a}. Its central abstraction is the \emph{Agent Card}, a machine-readable metadata document that each agent publishes to advertise its capabilities. An Agent Card contains structured fields describing the agent's identity, natural-language descriptions of its competencies, supported interaction modalities, and authentication requirements. In contrast to MCP tool definitions, which specify fine-grained callable functions and their parameter schemas, Agent Cards characterize coarse-grained, goal-level capabilities of an entire agent without prescribing its internal implementation. Consuming agents use these cards to identify suitable peers, negotiate task delegation through structured message exchanges, and coordinate multi-step workflows via shared artifacts. Importantly, A2A is explicitly agnostic to the internal architecture of participating agents: An A2A-registered agent may rely on MCP for its own tool access while exposing only high-level capabilities through its Agent Card. The two protocols thus occupy distinct but complementary roles, where MCP governs the \emph{tool plane}, defining how individual functions are described and invoked, whereas A2A governs the \emph{collaboration plane}, defining how agents advertise, discover, and delegate work to one another~\cite{ehtesham2025surveyagentinteroperabilityprotocols}.